\chapter{Estado actual y objetivos iniciales}
\chaptermark{Estado actual y objetivos}

\section{Motivación y antecedentes}
Este proyecto surge como una ampliación del trabajo desarrollado durante las
prácticas de empresa realizadas en el \gls{itc}. La librería
``\textbf{\gls{ngtgl}}'' fue diseñada originalmente para facilitar y acelerar
el desarrollo \gls{frontend} sobre un \gls{backend} específico, el del proyecto
\gls{nextgendem}, dedicado a la gestión de datos bioinformáticos.

Este planteamiento limitaba su aplicabilidad a proyectos con
arquitecturas prácticamente idénticas: cualquier modificación en el backend
obligaba a reescribir una cantidad notable de código en el frontend, porque
los componentes conocían de forma explícita la forma de los datos y los
endpoints concretos del proyecto original. Reconociendo esta limitación, y
tras consultarlo con mi tutor de empresa del \gls{itc}, se planteó la posibilidad de ampliar el
alcance del proyecto para construir un \gls{framework} completo que incluyese
tanto backend como frontend, de manera que pudiera reutilizarse en una
variedad mucho más amplia de proyectos y arquitecturas. Además, la llegada
de un nuevo proyecto con el \gls{jardincanario} manifestó la utilidad de un proyecto
de este tipo.

Por otra parte, el desarrollo ya estaba avanzado y, debido a las limitaciones
temporales de las prácticas de empresa, tuve que redirigir el trabajo hacia un
punto intermedio, ya que prioricé entregar la librería \gls{ngtgl} funcional. Así, el
proyecto original se orientó a crear una librería frontend de componentes reutilizables. 
Sobre la base de ese desarrollo se planteó el presente Trabajo de Fin de Grado, con el objetivo de consolidar,
documentar y presentar el framework de forma completa, backend y frontend, y demostrar su utilidad ante un modelo de datos en evolución.

\subsection{El ecosistema previo: NextGenDem, uniback y ngt-gui}
El punto de partida lo forman dos artefactos. Por un lado, un backend heredado
en Python basado en el backend de NextGenDem con una base de modelos de dominio (objetos funcionales,
autenticación y autorización, anotaciones, tareas) pero sin una capa de
comunicación uniforme ni mecanismos de exposición automática. Por otro, la
librería Angular \gls{ngtgl}, con un catálogo amplio de componentes de
presentación (tablas, menús, diálogos, layout) pero acoplada a un servicio de
backend monolítico con más de doscientos métodos escritos a mano, uno por cada
operación concreta del proyecto original.

\subsection{La colaboración con el Jardín Botánico Viera y Clavijo}
El \gls{jbvc} es la institución de conservación e investigación botánica principal
 en Canarias. Sus necesidades --- catalogación de especies, gestión
de colecciones, anotaciones de campo, jerarquías taxonómicas --- son un caso
de uso ideal de un modelo de datos en evolución ya que el esquema crece y se ajusta
a medida que avanza la digitalización de sus fondos y los propios estándares en los que se
definen los datos. El framework no es la
solución final para el Jardín, sino la herramienta con la que dicha solución
se construye; por ello la interacción con esta institución sirvió para calibrar
el grado de abstracción y genericidad que el framework debía ofrecer.

\section{Estado del desarrollo}
El desarrollo del framework se encontraba, al inicio del TFG, en una fase
intermedia y \emph{asimétrica}. El frontend tenía una buena parte de sus
funcionalidades implementadas, incluyendo componentes básicos y algunas
funcionalidades avanzadas, pero seguía acoplado al backend original. El
backend, que no se había abstraído, funcionaba como un backend tradicional: 
se me entregó una base
razonablemente completa de modelos ORM, pero sin una capa de API homogénea,
porque el \gls{itc} había priorizado disponer de un backend mínimamente
funcional para arrancar su último proyecto con recursos limitados.

En resumen, el framework no estaba terminado y, a excepción del frontend,
no disponía de una base de código suficiente para considerarse utilizable de forma genérica. 
Es aquí donde comienza mi
desarrollo en el presente Trabajo de Fin de Grado, con el objetivo de
completar, optimizar y cerrar tanto el backend como el frontend, y de unirlos mediante un contrato que los desacople del dominio. La figura~\ref{fig:fases-desarrollo} resume esta evolución.

\begin{figure}[!htbp]
\centering
\begin{tikzpicture}[
  node distance=7mm,
  fase/.style={rectangle, rounded corners, draw=black, thick,
    text width=11.5cm, align=left, inner sep=8pt, fill=gray!8},
  flecha/.style={-{Stealth[length=3mm]}, thick}
]
\node[fase] (partida) {{\bfseries Punto de partida (lo que había)}\\[2pt]
  Backend heredado en Python (modelos ORM por dominio, autenticación, anotaciones y tareas) \emph{sin} una capa de comunicación uniforme ni exposición automática.\\[2pt]
  Librería Angular \emph{ngt-gui} con un catálogo amplio de componentes, pero acoplada a un backend monolítico con más de doscientos métodos escritos a mano.};
\node[fase, below=of partida] (practicas) {{\bfseries Prácticas en el ITC (previo al TFG)}\\[2pt]
  Desarrollo del frontend \emph{ngt-gui}: catálogo de componentes reutilizables extraídos del proyecto NextGenDem y entrega de una librería frontend funcional.};
\node[fase, below=of practicas] (tfg) {{\bfseries Trabajo de Fin de Grado (este trabajo)}\\[2pt]
  API modular (envoltorio de respuesta uniforme, catálogo de errores y fábricas CRUD), contrato back--front con descubrimiento de entidades y capa cliente genérica, generación dinámica de formularios, generación automática de esquemas desde el ORM y desacople del dominio.};
\draw[flecha] (partida) -- (practicas);
\draw[flecha] (practicas) -- (tfg);
\end{tikzpicture}
\caption{Evolución del framework: punto de partida, desarrollo durante las prácticas en el ITC y aportación del presente Trabajo de Fin de Grado.}
\label{fig:fases-desarrollo}
\end{figure}

\section{Planificación previa al desarrollo}
Antes de iniciar el desarrollo se realizó una planificación preliminar con el
objetivo de establecer una hoja de ruta clara y estructurada que permitiera,
posteriormente, una planificación completa del proyecto.

\subsection{Análisis previo}

\subsubsection{Frontend}
Dado que el desarrollo del frontend se realizó durante las prácticas de
empresa, se conocía con detalle su alcance, su arquitectura interna y las
tareas necesarias para su finalización, lo que redujo la incertidumbre de
estimación de esta parte.

\subsubsection{Backend}
El backend, al ser un desarrollo que no realicé yo, presentaba mayor
incertidumbre: no tenía un conocimiento profundo de su alcance ni de las
tareas pendientes. Por ello contacté con mi tutor de empresa en el \gls{itc}
para obtener una visión clara del estado real, de la deuda técnica y de los
objetivos prioritarios.

\subsection{Acuerdo de planificación evolutiva}
Tras el contacto con la empresa se estableció una metodología iterativa incremental: 
a medida que avanzase el desarrollo y surgiesen nuevas necesidades
o desafíos, se irían ajustando y refinando los objetivos y tareas. Esta
planificación dinámica es coherente con la naturaleza del problema --- un
modelo de datos que evoluciona --- y permite mantener el proyecto alineado con
las necesidades reales en cada momento.

\subsection{Tareas iniciales}
Se definieron las tareas iniciales a partir del análisis previo y del
\emph{feedback} del \gls{itc}. Debido a la naturaleza evolutiva de la
planificación, estas tareas llevaron a divergir ligeramente, en ocasiones, de la planificación inicial
propuesta en el TFT01 y a elaborar una planificación más ajustada a las
necesidades reales, que se detalla en el capítulo de Desarrollo.

\section{Objetivos}
Tras contactar con mi tutor de prácticas, mi tutor académico y el
\gls{jardincanario}, se establecieron los objetivos generales siguientes:
\begin{enumerate}
    \item Diseñar e implementar un framework backend/frontend reutilizable.
    \item Proporcionar una arquitectura tecnológica extensible y desacoplada
    del dominio.
    \item Adaptar dinámicamente el desarrollo para asegurar su utilidad ante
    necesidades emergentes.
\end{enumerate}
Asimismo se definieron los siguientes objetivos específicos:
\begin{enumerate}
    \item \textbf{Crear una API modular.} Establecer un envoltorio de
    respuesta uniforme, un catálogo de errores estándar y fábricas de
    endpoints CRUD que eliminen el código repetitivo por entidad.
    \item \textbf{Diseñar y modelar los esquemas de datos.} Consolidar el
    núcleo de modelos ORM y dotarlos de versionado e histórico.
    \item \textbf{Implementar un sistema de comunicación estandarizado} entre
    backend y frontend, incluyendo descubrimiento de entidades y una capa
    cliente genérica tipada.
    \item \textbf{Diseñar, refactorizar y mejorar el generador automático de
    formularios} basado en los componentes Formly y Grid, guiado por el
    esquema publicado por el backend.
    \item \textbf{Refactorizar o reimplementar y mejorar los mecanismos para
    la generación automática de esquemas} a partir del modelo ORM, de forma
    que el modelo de datos sea la única fuente de verdad.
\end{enumerate}
